%% ===========================================================================
%%  4. Experiments
%%  Quantifying with Confidence: Conformal Prediction for Reliable Musical
%%  Style Similarity Ranking
%% ===========================================================================

\section{Experiments}
\label{sec:experiments}

We evaluate the proposed framework along three dimensions:
(1)~the quality of MERT embeddings for capturing stylistic similarity,
(2)~the coverage and set size characteristics of conformal prediction sets,
and (3)~the comparison of CP against bootstrap and Gaussian approximation
baselines. All experiments use the GTZAN dataset under a stratified 5-fold
cross-validation protocol.

\subsection{Dataset and Preprocessing}

\textbf{GTZAN}~\cite{tzanetakis2002gtzan} is the de facto standard benchmark
for music genre classification and a common testbed for music information
retrieval research. It consists of 1{,}000~audio tracks evenly distributed
across 10~genres: blues, classical, country, disco, hip-hop, jazz, metal,
pop, reggae, and rock. Each track is a 30-second monophonic WAV file sampled
at 22.05\,kHz.

During data validation, we identified one corrupted file
(\texttt{jazz.00054.wav}) that could not be successfully decoded, leaving
$N = 999$ tracks for the experiments. The remaining jazz class has 99 tracks;
all other classes have 100~tracks. This minor imbalance is accounted for by
our stratified splitting procedure and is negligible for the coverage analysis.

\textbf{Preprocessing.} All audio tracks are resampled to MERT's native
sampling rate of 24\,kHz, converted to mono via channel averaging, and
processed through the embedding extraction pipeline described in
Section~\ref{sec:method}. The output is a matrix
$\mathbf{E} \in \mathbb{R}^{999 \times 1024}$ of $\ell_2$-normalized
embeddings. Extraction of all embeddings requires approximately 9~minutes on
an NVIDIA RTX~5070~GPU (8\,GB~VRAM).

\subsection{Experimental Setup}

We adopt a stratified 5-fold protocol that preserves per-genre proportions
across all subsets. For each fold, the data is partitioned into three disjoint
splits:
\begin{itemize}[leftmargin=*,itemsep=1pt]
    \item \textbf{Reference set} ($60\%$): used to compute per-genre prototype
    centroids $\boldsymbol{\mu}_k$ via Equation~\ref{eq:centroid}.
    \item \textbf{Calibration set} ($20\%$): used to compute the empirical
    distribution of nonconformity scores $\mathcal{A}_{\text{cal}}$ and
    calibrate prediction thresholds.
    \item \textbf{Test set} ($20\%$): held out for evaluation; never used
    during centroid construction or calibration.
\end{itemize}
For each test point $x$, we compute nonconformity scores
$A(\mathbf{e}_x, k)$ with respect to all $K = 10$ genre centroids, derive
$p$-values via Equation~\ref{eq:pvalue}, and construct prediction sets
$\Gamma_{\alpha}(x)$ via Equation~\ref{eq:prediction_set}. Results are
aggregated across all five folds to produce means and standard deviations.

We evaluate at three significance levels $\alpha \in \{0.01, 0.05, 0.10\}$,
corresponding to confidence targets of $99\%$, $95\%$, and $90\%$,
respectively. Random seeds are fixed (base seed 42, incremented by fold
index) to ensure reproducibility.

\subsection{Embedding Quality Verification}

Before evaluating the conformal prediction framework, we assess whether
MERT embeddings capture genre-discriminative structure in the GTZAN dataset.
We compute two quantities: (i)~\emph{intra-class cosine similarity}---the
mean pairwise cosine similarity between embeddings of tracks belonging to
the same genre; and (ii)~\emph{inter-class cosine similarity}---the mean
cosine similarity between embeddings of tracks belonging to different genres.

The intra-class similarity averaged across all 10 genres is $0.917$, while
the inter-class similarity is $0.873$, yielding a separation margin of
$\Delta = 0.044$. The positive margin confirms that tracks from the same
genre are systematically closer in the MERT embedding space than tracks from
different genres. A two-dimensional PCA projection of the embedding matrix
(see Figure~\ref{fig:pca} in Appendix~\ref{sec:appendix}) further reveals
clear clustering tendencies, with genres such as \emph{classical} and
\emph{metal} forming tight, well-separated groups and more ambiguous genres
such as \emph{rock} and \emph{country} occupying overlapping regions---a
pattern consistent with known perceptual ambiguities in musical genre
taxonomy.

We detect no NaN or infinite values in the embedding matrix, confirming that
the MERT extraction pipeline is numerically stable.

\subsection{Coverage Analysis}

Table~\ref{tab:coverage_results} reports the empirical coverage of CP,
Bootstrap, and Gaussian methods averaged over the five cross-validation folds.

\begin{table}[htbp]
\centering
\caption{Empirical coverage and average prediction set size across three
significance levels (\textbf{bold} = best coverage). Coverage $\geq$ nominal
is highlighted as matching or exceeding the $1-\alpha$ target. Results
aggregated over 5 stratified folds (mean $\pm$ std).
}
\label{tab:coverage_results}
\begin{tabular}{c c c c c}
\toprule
$\alpha$ & Nominal $1-\alpha$ & CP & Bootstrap & Gaussian \\
\midrule
\addlinespace[2pt]
\multirow{2}{*}{0.01} & \multirow{2}{*}{0.99}
    & \textbf{0.993 $\pm$ 0.010} & 0.987 $\pm$ 0.012 & 0.965 $\pm$ 0.020 \\
    & & \hspace{2pt}(set size: 8.51) & (8.11) & (7.05) \\[4pt]
\multirow{2}{*}{0.05} & \multirow{2}{*}{0.95}
    & \textbf{0.955 $\pm$ 0.021} & 0.950 $\pm$ 0.023 & 0.932 $\pm$ 0.033 \\
    & & \hspace{2pt}(set size: 6.69) & (6.59) & (6.03) \\[4pt]
\multirow{2}{*}{0.10} & \multirow{2}{*}{0.90}
    & \textbf{0.910 $\pm$ 0.028} & \textbf{0.910 $\pm$ 0.029} & 0.907 $\pm$ 0.031 \\
    & & \hspace{2pt}(set size: 5.40) & (5.39) & (5.35) \\
\bottomrule
\end{tabular}
\end{table}

\textbf{Conformal prediction meets its coverage guarantee.} At
$\alpha = 0.01$, CP achieves coverage of $0.993$, exceeding the nominal
$0.99$ target. At $\alpha = 0.05$, coverage is $0.955$, and at $\alpha =
0.10$, it is $0.910$---both tightly matching their respective nominal levels.
The standard deviations (0.010--0.028) reflect the fold-to-fold variability
expected from a 999-track dataset, and are consistent with the
$\mathcal{O}(1/\sqrt{n})$ scaling of empirical coverage estimates.

\textbf{Gaussian approximation fails at high confidence.} The Gaussian
baseline achieves only $0.965$ coverage at $\alpha = 0.01$---a shortfall of
$2.5$~percentage~points below the $0.99$ nominal target. This confirms that
the nonconformity score distribution in the MERT embedding space deviates
from normality, exhibiting heavier tails than a Gaussian distribution. At
more lenient significance levels ($\alpha = 0.05, 0.10$), the Gaussian
approximation approaches acceptable coverage, consistent with the asymptotic
robustness of Gaussian quantiles in the central regime. However, the
deficiency at $\alpha = 0.01$ is precisely the regime most relevant for
high-stakes filtering and safety-critical applications.

\textbf{Bootstrap performs comparably to CP but without guarantees.} The
bootstrap baseline achieves coverage of $0.987$ at $\alpha = 0.01$, slightly
below the nominal target, and $0.950$ and $0.910$ at $\alpha = 0.05$ and
$\alpha = 0.10$, respectively. While its aggregate performance closely
tracks CP on the GTZAN dataset, bootstrap offers no finite-sample coverage
guarantee and requires $B = 1{,}000$ resamples of the calibration set,
incurring roughly $1{,}000\times$ the computational cost of CP calibration.

\subsection{Prediction Set Size Analysis}

The second row for each $\alpha$ in Table~\ref{tab:coverage_results} reports
the average prediction set size. The key observations are:

\begin{itemize}[leftmargin=*,itemsep=1pt]
    \item \textbf{Set size decreases with $\alpha$, as expected.} As the
    significance level relaxes ($\alpha: 0.01 \to 0.10$), the
    $(1 - \alpha)$-quantile of calibration scores decreases, admitting fewer
    genres into the prediction set. CP set sizes shrink from $8.51$ (at
    $\alpha = 0.01$) to $5.40$ (at $\alpha = 0.10$), reflecting increasing
    permissiveness of the inclusion criterion.

    \item \textbf{Gaussian achieves smaller sets through undercoverage.} The
    Gaussian baseline consistently produces the smallest prediction sets
    (e.g., $7.05$ vs.\ CP's $8.51$ at $\alpha = 0.01$), but this apparent
    advantage is spurious: it stems from an overly optimistic threshold
    derived from the normality assumption, which excludes true labels that
    fall in the heavy tail. The smaller set size is not a sign of higher
    precision but of \emph{overconfidence}---a well-known pitfall of
    parametric methods on non-Gaussian data.

    \item \textbf{CP achieves the best coverage--size trade-off.} At
    $\alpha = 0.01$, CP's set size of $8.51$ is approximately $1.46$ larger
    than the Gaussian baseline's $7.05$, yet CP provides valid coverage
    ($0.993 \geq 0.99$) while Gaussian does not ($0.965 < 0.99$). This
    demonstrates that CP pays an appropriate ``uncertainty premium'' in set
    size to deliver its distribution-free coverage guarantee---a premium that
    is both principled and empirically necessary.
\end{itemize}

The distribution of set sizes across individual test tracks (visualized in
Figure~\ref{fig:set_size_boxplot} in Appendix~\ref{sec:appendix}) reveals
substantial heterogeneity: while most tracks receive prediction sets
containing 4--9 genres, a small number of tracks are assigned either very
small sets (1--2 genres, indicating high confidence) or the full set of 10
genres (indicating maximal ambiguity). Empty prediction sets are rare
($<1\%$ across all folds and $\alpha$ levels), consistent with the fact that
GTZAN tracks are drawn from the same distribution as the calibration data.

\subsection{Cross-Dataset Validation on FMA}
\label{sec:fma_experiment}

To assess whether the conformal prediction framework generalizes beyond
GTZAN, we conduct a parallel experiment on the Free Music Archive (FMA) Small
dataset~\cite{defferrard2017fma}---a substantially larger and more diverse
collection of full-length Creative Commons-licensed music. FMA Small contains
8{,}000 tracks (7{,}994 after excluding corrupted files) spanning eight
top-level genres: Electronic, Experimental, Folk, Hip-Hop, Instrumental,
International, Pop, and Rock. Tracks are variable-length MP3 files (most 30--180\,s)
encoding a broader range of musical styles and recording conditions than
GTZAN's curated 30-second excerpts.

\textbf{Embedding quality.} Following the same MERT extraction pipeline
(Section~\ref{sec:method}), we obtain a $7{,}994 \times 1024$ embedding matrix
for FMA Small. The intra-class cosine similarity is $0.847$, the inter-class
similarity is $0.821$, yielding a separation margin of $\Delta = 0.026$. This
is substantially narrower than GTZAN's $\Delta = 0.044$---approximately
$41\%$ weaker separation---consistent with the larger, noisier, and
stylistically more heterogeneous composition of the FMA dataset. The positive
margin nonetheless confirms that MERT embeddings exhibit meaningful
genre-discriminative structure.

\textbf{CP results on FMA.} Table~\ref{tab:fma_results} reports the empirical
coverage and set size for all three methods on FMA, aggregated over five
stratified folds.

\begin{table}[htbp]
\centering
\caption{Empirical coverage and average prediction set size on the FMA Small
dataset (\textbf{bold} = best coverage). Results aggregated over 5 stratified
folds (mean $\pm$ std). FMA contains more tracks ($\sim$8{,}000) and exhibits
weaker class separation ($\Delta = 0.026$) than GTZAN.
}
\label{tab:fma_results}
\begin{tabular}{c c c c c}
\toprule
$\alpha$ & Nominal $1-\alpha$ & CP & Bootstrap & Gaussian \\
\midrule
\addlinespace[2pt]
\multirow{2}{*}{0.01} & \multirow{2}{*}{0.99}
    & \textbf{0.988 $\pm$ 0.004} & 0.987 $\pm$ 0.004 & 0.966 $\pm$ 0.003 \\
    & & \hspace{2pt}(set size: 7.75) & (7.72) & (7.22) \\[4pt]
\multirow{2}{*}{0.05} & \multirow{2}{*}{0.95}
    & \textbf{0.948 $\pm$ 0.004} & 0.947 $\pm$ 0.004 & 0.934 $\pm$ 0.005 \\
    & & \hspace{2pt}(set size: 6.85) & (6.83) & (6.59) \\[4pt]
\multirow{2}{*}{0.10} & \multirow{2}{*}{0.90}
    & 0.898 $\pm$ 0.006 & 0.896 $\pm$ 0.006 & \textbf{0.903 $\pm$ 0.005} \\
    & & \hspace{2pt}(set size: 6.03) & (6.02) & (6.10) \\
\bottomrule
\end{tabular}
\end{table}

Several observations are noteworthy. First, CP continues to provide coverage
closely matching nominal targets despite the weaker embedding space---$0.988$
at $\alpha = 0.01$ (target $0.99$) and $0.948$ at $\alpha = 0.05$ (target
$0.95$). This demonstrates that the distribution-free guarantee is robust to
the quality of the underlying representation: CP adapts its threshold to the
empirical score distribution regardless of the discriminability of the
embedding space.

Second, the standard deviations on FMA ($\sim$0.003--0.006) are substantially
smaller than on GTZAN ($\sim$0.010--0.033), reflecting the nearly $8\times$
larger sample size and the corresponding reduction in fold-to-fold variance.

Third, the Gaussian baseline at $\alpha = 0.10$ achieves slightly higher
coverage ($0.903$) than CP ($0.898$), an inversion of the pattern observed on
GTZAN. This occurs because the Gaussian quantile is more conservative than CP's
empirical quantile at this particular significance level for the FMA score
distribution. However, at $\alpha = 0.01$---the high-confidence regime most
critical for downstream applications---Gaussian undercovers by $2.4$ percentage
points ($0.966$ vs.\ target $0.99$), consistently failing where guarantees
matter most.

Finally, the prediction set sizes on FMA are generally smaller than on GTZAN
(e.g., $7.75$ vs.\ $8.51$ at $\alpha = 0.01$), despite the weaker class
separation. This counterintuitive result can be explained by two factors:
(1)~the substantially larger calibration set ($\sim$1{,}600 tracks vs.\
$\sim$200 on GTZAN) provides a more precise empirical quantile estimate,
enabling tighter thresholds; and (2)~FMA's eight genres present a simpler
prediction task ($K = 8$) than GTZAN's ten ($K = 10$), reducing the
theoretical upper bound on set size.

\subsection{Summary of Experimental Findings}

Our experiments yield four primary takeaways:

\begin{enumerate}[leftmargin=*,itemsep=1pt]
    \item \textbf{MERT embeddings are genre-discriminative.} On GTZAN, the
    intra-class similarity ($0.917$) meaningfully exceeds the inter-class
    similarity ($0.873$), confirming that MERT captures stylistic structure.
    On FMA, the weaker separation ($\Delta = 0.026$) reflects the dataset's
    greater diversity but remains above chance, with CP adapting to the
    reduced discriminability without compromising its coverage guarantee.

    \item \textbf{Conformal prediction provides valid, tight coverage across
    datasets.} At all three significance levels on both GTZAN and FMA, CP
    achieves empirical coverage at or near the nominal target, with only
    moderate set sizes. On FMA, CP achieves $0.988$ coverage at $\alpha=0.01$
    despite the substantially weaker embedding space ($\Delta=0.026$ vs.\
    GTZAN's $\Delta=0.044$), demonstrating that the distribution-free
    guarantee is representation-agnostic.

    \item \textbf{Parametric baselines fail where CP succeeds.} The Gaussian
    approximation baseline undercovers at $\alpha = 0.01$ on \emph{both}
    datasets---$0.965$ on GTZAN and $0.966$ on FMA (vs.\ target $0.99$).
    The pervasiveness of this failure across datasets with markedly different
    score distributions underscores that the normality assumption is
    fundamentally inappropriate for perceptual similarity tasks. Bootstrap
    resampling, while empirically competitive, lacks the finite-sample
    validity guarantee that distinguishes CP.

    \item \textbf{Larger calibration sets yield tighter prediction sets.} The
    FMA calibration set ($\sim$1{,}600 tracks) is approximately $8\times$
    larger than GTZAN's ($\sim$200 tracks), resulting in prediction sets that
    are consistently smaller on FMA ($7.75$ vs.\ $8.51$ at $\alpha=0.01$)
    despite weaker class separation. This illustrates a key practical
    principle: deploying CP with a sufficiently large calibration corpus
    narrows the gap between the coverage guarantee and practical
    informativeness.
\end{enumerate}
